% Faithful formal-publication transcription: main statement, Section2 and Richardson machinery.
\section*{Author main result}
\begin{authorclaim}{Theorem2 (main result, see also Corollary15)}
Suppose PRG is an $(n,w,n^{-2})$ PRG with seed length $s_0$, computable in space $m$. Then, for every $\epsilon$ there exists an $(n,w,\epsilon)$ WPRG with seed length
\[s=s_0+O\left(\log\frac w\epsilon\cdot\log\log_n\frac1\epsilon\right)\]
that is computable in space $O(m+(\log\log(w/\epsilon))^3)$.
\end{authorclaim}
\section*{Necessary ROBP and weighted-distribution context (Sections1.2--1.3)}
A length $n$, width $w$ read-once branching program (ROBP) is a directed graph whose nodes, called states, are partitioned to $n+1$ layers, each consists of at most $w$ states. The first layer contains a designated ``start'' state, and the last layer consists of two states labeled `accept' and `reject'. From every state but for the latter two, there are two outgoing edges, labeled by 0 and 1, to the following layer.\footnote{For simplicity, here we only consider ROBPs with two outgoing edges. Larger out-degrees (or alphabet) can also be considered and is in fact crucial for obtaining our result even if one is only interested in the binary case.} On input $x\in\{0,1\}^n$, the computation proceeds by following the edges according to the labels given by the bits of $x$ starting from the start state. The string $x$ is accepted by the program if the computation ends in the accept state.
We say that a distribution $D$ on $n$-bit strings is $(n,w,\epsilon)$ pseudorandom if for every length $n$, width $w$ ROBP, the path induced by an instruction sequence that is sampled from $D$ has, up to an additive error $\epsilon$, the same probability to end in the accept state as a truly random path. A truly random path corresponds to a path picked uniformly at random from the $2^n$ possible paths. An $(n,w,\epsilon)$ pseudorandom generator (PRG) is an algorithm $\mathrm{PRG}:\{0,1\}^s\to\{0,1\}^n$ that when fed with $s$ uniformly random bits has an output distribution that is $(n,w,\epsilon)$ pseudorandom. We refer to the input to PRG as the seed.
\begin{authorclaim}{Definition1 (pseudorandom pseudo-distribution)}
Let $\rho_1,\ldots,\rho_{2^s}\in\mathbb R$ and $p_1,\ldots,p_{2^s}\in\{0,1\}^n$. The sequence $\widetilde D=((\rho_1,p_1),\ldots,(\rho_{2^s},p_{2^s}))$ is an $(n,w,\epsilon)$ pseudorandom pseudo-distribution (PRPD) if for every length $n$, width $w$ ROBP, the sum of all $\rho_i$-s for which the respective paths $p_i$ end in the accept state is an $\epsilon$-approximation to the probability of ending in the accept state by taking a truly random path in the program.
\end{authorclaim}
Note that Definition1 allows the weights $\rho_i$ to take both positive and negative values. These values are not necessarily bounded by 1 in absolute value, nor by any constant for that matter, and they do not necessarily sum up to 1. Nevertheless, the definition requires that the numbers cancel out nicely so that summing the weights of the respective paths that arrive to the accept state yields an $\epsilon$-approximation for the probability of arriving to the accept state by taking a truly random path (and, in particular, the sum is a number in $[-\epsilon,1+\epsilon]$). Analogous to a PRG, an $(n,w,\epsilon)$ weighted pseudorandom generator (WRPG) is an algorithm $\mathrm{WPRG}:\{0,1\}^s\to\mathbb R\times\{0,1\}^n$ whose output, when fed with a uniform seed, is an $(n,w,\epsilon)$ PRPD.
A WPRG that can be computed in bounded space suffices to derandomize two-sided error randomized algorithms. Indeed, the straightforward derandomization using a pseudorandom (proper) distribution, which sums the probability mass of the relevant paths, works just as well for pseudo-distributions as one can sum up the weights $\rho_i$ which, in a sense, generalize the probability mass. Of course, the space requirement now depends on the bit complexity of the weights as well.
\section*{2 Preliminaries}
\subsection*{2.1 Matrices, branching programs, and space complexity}
A matrix is Boolean if all its entries are in $\{0,1\}$, and stochastic if all its entries are nonnegative and the sum of each column is 1. Denote by $\operatorname{BSto}(w)$ the set of $w\times w$ boolean stochastic matrices. We will denote by $\|\cdot\|$ the induced $\ell_1$ norm, i.e., $\|A\|=\max_j\sum_i|A_{i,j}|$.
We will often work with block matrices. For instance, we may interpret $A\in\mathbb R^{nm\times nm}$ as an $n\times n$ matrix with entries which are $m\times m$ matrices. Whenever this interpretation is clear, we let $A[i,j]$ be the $(i,j)$-th block. In this example, $A[i,j]\in\mathbb R^{m\times m}$.
\begin{authorclaim}{Definition4 (branching program)}
Let $\Sigma$ be some alphabet and let $n,w\in\mathbb N$. An $(n,\Sigma,w)$ branching program (BP) is a sequence $B=(B_1,\ldots,B_n)$, where each $B_i:\Sigma\to\operatorname{BSto}(w)$.
\end{authorclaim}
For $b\le a$ we let $B_{[b,a]}$ be the $(a-b+1,\Sigma,w)$ BP $(B_a,\ldots,B_b)$.
\begin{authorclaim}{Definition5}
The value of an $(n,\Sigma,w)$ BP $B=(B_1,\ldots,B_n)$ on $x=(x_1,\ldots,x_n)\in\Sigma^n$, denoted $\operatorname{val}(B,x)$, is the realized $w\times w$ matrix of $B$ when fed by $x$, i.e.
\[\operatorname{val}(B,x)=B_n(x_n)\cdot B_{n-1}(x_{n-1})\cdots B_1(x_1).\]
If $B$ is the empty sequence, we set $\operatorname{val}(\emptyset,x)=I_w$.
\end{authorclaim}
\begin{authorclaim}{Definition6 (weighted PRG)}
We say $W$ is an $(n,\Sigma,w,\epsilon)$-WPRG against BPs with seed length $s$ if:
\begin{itemize}
\item $W=(I,\mu)$ where $I:\{0,1\}^s\to\Sigma^n$ and $\mu:\{0,1\}^s\to\mathbb R$, and,
\item For every $(n,\Sigma,w)$ BP $B=(B_1,\ldots,B_n)$, it holds that
\[\left\|\mathbb E_{x\in\{0,1\}^s}[\mu(x)\cdot\operatorname{val}(B,I(x))]-\mathbb E_{x\in\Sigma^n}[\operatorname{val}(B,x)]\right\|\le\epsilon.\]
\end{itemize}
\end{authorclaim}
When $\mu\equiv1$, we say that $W$ is a PRG.
For $1\le\ell\le n$ we let $G_\ell:\{0,1\}^{s_0}\to\Sigma^\ell$ be the first $\ell$ symbols of the output of $G$. Note that if $G:\{0,1\}^{s_0}\to\Sigma^n$ is an $(n,\Sigma,w,\epsilon)$ PRG then $G_\ell$ is an $(\ell,\Sigma,w,\epsilon)$ PRG.
We say $f:\Lambda_1\to\Lambda_2$ is computable in space $s$, if given $x\in\Lambda_1$ and index $j$, $f(x)_j\in\Lambda_2$ can be computed in additional work space that consists of $s$ bits. We will use the following well known theorem regarding the space complexity of compositions.
\begin{authorclaim}{Theorem7}
Let $f_1,f_2:\{0,1\}^*\to\{0,1\}^*$ be two functions that can be computed in $s_1,s_2:\mathbb N\to\mathbb N$ space such that $s_1(n),s_2(n)=\Omega(\log n)$. Then, on input $x$, $f_2\circ f_1:\{0,1\}^*\to\{0,1\}^*$ can be computed using $O(s_1(|x|)+s_2(|f_1(x)|))$ space.
\end{authorclaim}
\subsection*{2.2 Known PRG constructions}
\begin{authorclaim}{Theorem8 ([17,18])}
For any positive integers $n,w$, any error parameter $\epsilon>0$ and any alphabet $\Sigma$, there exists an $(n,\Sigma,w,\epsilon)$ PRG with seed length
\[s=O\left(\log n\cdot\log\frac{nw|\Sigma|}\epsilon\right),\]
computable in space $\min\{O(\log(nw|\Sigma|/\epsilon)),O(\log n\cdot\log\log(nw|\Sigma|/\epsilon))\}$.
\end{authorclaim}
\begin{authorclaim}{Theorem9 ([15])}
For any positive integers $n,w$, any error parameter $\epsilon>0$ and any alphabet $\Sigma$, there exists an $(n,\Sigma,w,\epsilon)$ PRG with seed length
\[s=\log|\Sigma|+O(\log n\cdot\log(nw/\epsilon)),\]
computable in space $O(\log n\cdot(\log\log(nw|\Sigma|/\epsilon))^2)$.
\end{authorclaim}
Theorem8 is derived almost directly from [17,18], and Theorem9 follows from [15], except for the space complexity which is implicit in those works and also depends on the specific implementation. For completeness, we give the proof of Theorem8 in AppendixB.1, and of Theorem9 in AppendixB.3.
\section*{3 Richardson iteration}
Let $A$ be an invertible $n\times n$ real matrix, and assume that $B$ approximates $A^{-1}$, concretely, $\|B-A^{-1}\|\le\epsilon_0$ for some sub-multiplicative norm. Richardson iteration is a method for obtaining a more refined approximation of $A^{-1}$ given access to the crude $B$ as well as to the original matrix $A$.
\begin{authorclaim}{Lemma10}
Let $L\in\mathbb R^{m\times m}$ be an invertible matrix and $A\in\mathbb R^{m\times m}$ such that $\|L^{-1}-A\|\le\epsilon_0$. For any nonnegative integer $k$, define
\[R(A,L,k)=\sum_{i=0}^k(I-AL)^iA.\]
Then, $\|L^{-1}-R(A,L,k)\|\le\|L^{-1}\|\cdot\|L\|^{k+1}\cdot\epsilon_0^{k+1}$.
\end{authorclaim}
The proof is deferred to AppendixA.
Following [1] we will be interested in the following instantiation of the Richardson iteration. Let $\mathbf M=(M_1,\ldots,M_n)$ be a sequence of $w\times w$ matrices. We consider the $(n+1)w\times(n+1)w$ matrix
\begin{equation}\tag{4}M=\begin{pmatrix}0&0&\cdots&0&0\\M_1&0&\cdots&0&0\\0&M_2&\ddots&0&0\\\vdots&\vdots&\ddots&\vdots&\vdots\\0&0&\cdots&M_n&0\end{pmatrix}.\end{equation}
The Laplacian of $\mathbf M$ is $L=I_{(n+1)w}-M$, and we treat $L$ as an $(n+1)\times(n+1)$ block matrix. The following claim follows by a simple calculation.
\begin{authorclaim}{Claim11}
For $i,j\in[n+1]$, the $(i,j)$-th block of $L^{-1}$ is given by
\[L^{-1}[i,j]=\begin{cases}M_{i-1}\cdots M_j&i>j,\\I_w&i=j,\\0&i<j.\end{cases}\]
\end{authorclaim}
\paragraph{Richardson for branching programs.} Let $B=(B_1,\ldots,B_n)$ be an $(n,\Sigma,w)$ BP and let $M_i=\mathbb E_{\sigma\in\Sigma}[B_i(\sigma)]$ be the corresponding transition matrices. Thus, approximating the transition probabilities of $B$,
\[\mathbb E_{x\in\Sigma^n}[\operatorname{val}(B,x)]=M_n\cdots M_1,\]
amounts to approximating the lowest leftmost entry $L^{-1}[n+1,1]$.
\begin{authorclaim}{Claim12}
Let $B=(B_1,\ldots,B_n)$ be an $(n,\Sigma,w)$ BP. Set $M_i=\mathbb E_{\sigma\in\Sigma}[B_i(\sigma)]$ and $L$ as in Equation(4). Also, let $G:\{0,1\}^s\to\Sigma^n$ be an $(n,\Sigma,w,\epsilon_0)$ PRG and consider
\begin{equation}\tag{5}A[a,b]=\begin{cases}\mathbb E_{x\in\{0,1\}^s}[\operatorname{val}(B_{[b,a-1]},G_{a-b}(x))],&a\ge b,\\0&a<b.\end{cases}\end{equation}
Then, $\|L^{-1}-R(A,L,k)\|\le(n+1)\cdot(2\epsilon_0)^{k+1}$.
\end{authorclaim}
Let $A$ as in Equation(5) and write $R(A,L,k)=\sum_{i=0}^k\Delta^iA$ where $\Delta=I-AL$. Denote $A'=A-I$, i.e., $A'$ is the part of $A$ below the main diagonal. Then,
\[\Delta=I-AL=I-A(I-M)=(I-A)+AM=AM-A'.\]
In block notation, for $a,b\in[n+1]$, following Equation(4),
\[AM[a,b]=\sum_{i=1}^{n+1}A[a,i]M[i,b]=A[a,b+1]M[b+1,b]=A[a,b+1]\cdot M_b.\]
Thus,
\begin{equation}\tag{6}\Delta[a,b]=\begin{cases}A[a,b+1]\cdot M_b-A[a,b]&a>b,\\0&a\le b.\end{cases}\end{equation}
Going back to $R(A,L,k)$, for $a>b$ we have that
\begin{equation}\tag{7}R(A,L,k)[a,b]=A[a,b]+\sum_{i=1}^k\sum_{a>r_i>\cdots>r_1\ge b}\Delta[a,r_i]\cdot\Delta[r_i,r_{i-1}]\cdots\Delta[r_2,r_1]\cdot A[r_1,b].\end{equation}
If we further let $C_0[a,b]=A[a,b+1]\cdot M_b$ and $C_1[a,b]=A[a,b]$, then
\begin{equation}\tag{8}\begin{aligned}R(A,L,k)[a,b]=A[a,b]+\sum_{\substack{t\in\{0,1\}^i\\a>r_i>\cdots>r_1\ge b}}\sum_{t_1,\ldots,t_i\in\{0,1\}}(-1)^{t_1+\cdots+t_i}\cdot C_{t_i}[a,r_i]\cdots C_{t_1}[r_2,r_1]\cdot A[r_1,b].\end{aligned}\end{equation}
